% Scope: Build an operator-based account of reconstruction from least squares to regularization and model-based inference; restrict learned methods to physics, stability and validation.
% Status: architecture only; no chapter prose or completed problems yet.
% Prerequisites: Chapters 1, 6 and 14 through 15.
% Planned worked examples: Tikhonov solution, adjoint consistency and iterative convergence.
% Planned figures: Singular spectra, bias-variance curves and reconstruction residuals.
% Planned challenge problems: Prove nonuniqueness; compare reconstructions with matched data fidelity.
\chapter{Reconstruction and Inverse Problems}
\label{ch:16}

\section{The discrete forward model}
\label{sec:16:the-discrete-forward-model}
\subsection{Encoding operators and discretization}
\subsection{Adjoint versus inverse and inner-product tests}
\subsection{Noise whitening and likelihood}

\section{Linear reconstruction}
\label{sec:16:linear-reconstruction}
\subsection{Least squares, pseudoinverses and rank}
\subsection{Tikhonov regularization and singular-value filtering}
\subsection{Conditioning, resolution matrices and uncertainty}

\section{Sparsity and compressed sensing}
\label{sec:16:sparsity-and-compressed-sensing}
\subsection{Transform sparsity and incoherent measurements}
\subsection{L1 and total-variation objectives}
\subsection{Sampling assumptions and failure cases}

\section{Iterative algorithms}
\label{sec:16:iterative-algorithms}
\subsection{Conjugate gradients and normal equations}
\subsection{Proximal gradients and variable splitting}
\subsection{Stopping rules, preconditioning and data consistency}

\section{Low-rank and model-based reconstruction}
\label{sec:16:low-rank-and-model-based-reconstruction}
\subsection{Dynamic subspaces and low-rank penalties}
\subsection{Joint image and parameter estimation}
\subsection{Nonconvexity, initialization and identifiability}

\section{Learned reconstruction}
\label{sec:16:learned-reconstruction}
\subsection{Learned priors and unrolled optimization}
\subsection{Distribution shift, hallucination and uncertainty}
\subsection{Validation with residuals and task-based metrics}

\section{Worked examples}
\label{sec:16:worked-examples}
% Planned calculations: Tikhonov solution, adjoint consistency and iterative convergence.
\subsection{Derivation and quantitative calculation}
\subsection{Assumptions, limiting cases and scanner interpretation}

\section{Challenge problems}
\label{sec:16:challenge-problems}
% Problem design: Prove nonuniqueness; compare reconstructions with matched data fidelity.
% Use challengeproblem and stable labels prob:16:<topic>.
% Complete solutions belong in Appendix E, section for this chapter.
% Use \begin{solution}{prob:16:<topic>} ... \end{solution} there.
\subsection{Derivation and quantitative design}
\subsection{Model criticism and integrated reasoning}
